\section{A common squareclass descent}\label{sec:common-descent}

The sum of the products of opposite distances gives a useful one-parameter
family of covering curves.  The four square roots in this family determine
one point of the square.  In particular, the squareclass of their common
coordinate can be compared at different places.

\subsection{The four-root curve and its return map}

Let $t\in\Q\setminus\{0,1,-1\}$, and put
\[
 s=t-1,\qquad T=t+1,\qquad d=sT=t^2-1.
\]
Let $\mathcal U_t$ be the affine curve
\begin{equation}\label{eq:common-four-roots}
 P^2=\ell+T,\qquad C^2=\ell-s,\qquad
 M^2=\frac{s(T-\ell)}\ell,\qquad B^2=\frac{T(\ell+s)}\ell,
 \qquad \ell PCMB\ne0.
\end{equation}
We denote its smooth projective completion by $D_t$.  The order of the
four roots will be fixed throughout this section.

For a point whose distances from $(0,0),(1,0),(0,1),(1,1)$ are
$r_0,r_1,r_2,r_3$, respectively, the coordinates are
\begin{equation}\label{eq:common-from-distances}
\begin{gathered}
 t=r_0r_3+r_1r_2,\qquad P=r_0+r_3,\qquad M=r_3-r_0,\\
 B=r_1+r_2,\qquad C=r_2-r_1,\qquad \ell=P^2-t-1.
\end{gathered}
\end{equation}
To verify the forward construction and its range, put $u=1-2x$,
$v=1-2y$, $A=r_0r_3$, and $A'=r_1r_2$.
Expansion of the four squared-distance identities gives
\[
 A^2-(A')^2=-uv,\qquad
 (A+A')^2+(A-A')^2=1+(u^2+v^2)^2/4.
\]
Thus $t^2=(A+A')^2$ is the larger root of
\[
 F_t(Z)=Z^2-\bigl(1+(u^2+v^2)^2/4\bigr)Z+u^2v^2.
\]
The diagonal exclusion in Lemma~\ref{lem:ell-diagonal-exclusion} gives
$u^2\ne v^2$, and hence $F_t(1)=-(u^2-v^2)^2/4<0$.
For $a=u^2,b=v^2\in[0,1)$,
\[
 16F_t(5/4)=5-5a^2+6ab-5b^2>0.
\]
Indeed this polynomial is strictly concave on $[0,1]^2$, with vertex
values $5,0,0,1$, and its two zero vertices are excluded.
The axis of $F_t$ is below $1$, so $1<t<\sqrt5/2$.
Furthermore
\[
 PM=u+v,\quad BC=v-u,\quad
 P^2-C^2=B^2-M^2=2t,\quad P^2+M^2=u^2+v^2+2.
\]
These identities give $(P^2-T)(M^2+s)=d$ by expansion.
The first two equations of \eqref{eq:common-four-roots} follow from the
definition of $\ell$; the last two then follow from this product
identity and $B^2-M^2=2t$.  The diagonal exclusion also gives
$C,M\ne0$.  Since $s>0$ and $d>0$, the product identity first gives
$\ell>0$, and the equations for $C^2,M^2$ then give $s<\ell<T$.

Conversely, a tuple in \eqref{eq:common-four-roots} determines
\begin{equation}\label{eq:common-return}
\begin{gathered}
 u=\frac{PM-BC}{2},\qquad v=\frac{PM+BC}{2},\qquad
 x=\frac{1-u}{2},\qquad y=\frac{1-v}{2},\\
 r_0=\frac{P-M}{2},\quad r_3=\frac{P+M}{2},\quad
 r_1=\frac{B-C}{2},\quad r_2=\frac{B+C}{2}.
\end{gathered}
\end{equation}
These formulas give the four squared-distance identities over every field
of characteristic zero.  Indeed, the root equations imply
\[
 B^2-M^2=P^2-C^2=2t,\qquad
 P^2M^2+B^2C^2=2(P^2+M^2-2).
\]
Consequently $u^2+v^2=P^2+M^2-2=B^2+C^2-2$.  Since
$u+v=PM$ and $u-v=-BC$, substitution into
$\bigl((1-u)^2+(1-v)^2\bigr)/4$ and its three translates gives
$r_0^2,r_1^2,r_2^2,r_3^2$.  Moreover,
\[
 r_0r_3+r_1r_2=\frac{P^2-M^2+B^2-C^2}{4}=t.
\]
For $t>1$, the strict real return chamber is
\begin{equation}\label{eq:common-chamber}
 s<\ell<T,\qquad P,B>0,\qquad |PM|+|BC|<2.
\end{equation}
In this chamber $P>|M|$ and $B>|C|$, so all four returned distances are
positive.  The last inequality is equivalent to $|u|,|v|<1$, and hence to
$0<x,y<1$.  These identities also show directly that the transformation
\begin{equation}\label{eq:common-reflection}
 (\ell,P,C,M,B)\longmapsto(d/\ell,B,M,C,P)
\end{equation}
is the reflection $x\mapsto1-x$ on the square.

\begin{proposition}\label{prop:common-genus-boundary}
The curve $D_t$ is geometrically connected of genus nine.  The map
$\ell:D_t\to\PP^1$ has degree sixteen and is branched precisely at
\[
 0,\quad\infty,\quad s,\quad T,\quad-s,\quad-T.
\]
If $t>1$, none of the points above these six values is rational.
\end{proposition}
\begin{proof}
Over $\overline\Q(\ell)$, the four radicands in
\eqref{eq:common-four-roots} have independent squareclasses: their
valuations at $-T,s,T,-s$, respectively, each isolate one of the four
classes.  The extension therefore has degree sixteen.  At each of the
six displayed branch values the inertia group has order two.  For
example, the valuation vectors at $0$ and $\infty$, in the order
$(P,C,M,B)$, are $(0,0,1,1)$ and $(1,1,0,0)$ modulo two.  Thus
Riemann--Hurwitz gives
\[
 2g(D_t)-2=-32+6\cdot8=16.
\]
The following table records a regular function on the normalization and
its value at each omitted fibre:
\[
\begin{array}{c|rrrrrr}
 \ell&s&T&-s&-T&0&\infty\\\hline
 \text{regular square}&M^2&C^2&P^2&B^2&C^2&M^2\\
 \text{value}&2&2&2&2&-s&-s.
\end{array}
\]
The first four require $\sqrt2\in\Q$, and the last two require a square
root of the negative number $-s$.  This excludes every rational point
on the six fibres.
\end{proof}

\subsection{Prime support and the ramified unit conditions}

Write $[a]_K$ for the class of $a\in K^\times$ modulo squares; omit the
subscript when $K=\Q$.  For a nonzero rational number $d$, set
\[
 S(d)=\{p:\ p\text{ is odd and }v_p(d)>0\text{ is odd}\}.
\]

\begin{lemma}[Support of the common coordinate]\label{lem:common-support}
If $(\ell,P,C,M,B)\in\mathcal U_t(\Q)$, then
\[
 [\ell]\in\langle[-1],[2],[p]:p\in S(d)\rangle.
\]
For $t>1$, the sign is positive and the generator $[-1]$ can be omitted.
\end{lemma}
\begin{proof}
Fix an odd prime $p$ and put $m=v_p(t)$ and $n=v_p(\ell)$.
Suppose first that $m<0$, so $v_p(s)=v_p(T)=m$.
If $n>m$, the first and third equations of
\eqref{eq:common-four-roots} have valuations $m$ and $2m-n$;
both must be even, so $n$ is even.  If $n<m$, the fourth and first
have valuations $m$ and $n$, giving the same conclusion.  If $n=m$,
the reductions of $P^2/t$ and $C^2/t$ differ by $2$.
At least one is a unit; its square has valuation $m=n$, which is even.

Now let $m\ge0$.  If both $s$ and $T$ are units, then $n<0$ is even
by the first equation and $n>0$ is even by the third, whose valuation
is $-n$.
If $e=v_p(s)>0$, then $T$ is a unit.  Negative $n$ is even by the
first equation, and $0<n<e$ is even by the second.  For $n>e$,
the second and third equations have valuations $e$ and $e-n$, so
both $e$ and $n$ are even.  Hence odd $n$ can occur only at
$n=e$ with $e$ odd.  In that case $e=v_p(d)>0$.
If instead $e=v_p(T)>0$, the first equation handles $n<0$ and
$0<n<e$, while the first and fourth have valuations $e,e-n$ for
$n>e$.  Again odd $n$ forces $n=e=v_p(d)>0$ odd.

This accounts for every odd prime, including equal-valuation cases.
Squareclasses in $\Q$ are determined by sign and valuation parity.
Finally $C^2=\ell-s$ implies $\ell>s>0$ when $t>1$.
\end{proof}

\begin{lemma}[Ramified classes]\label{lem:common-ramified}
Let $p$ be odd and suppose $v_p(d)>0$ is odd.  If
$\mathcal U_t(\Q_p)$ is nonempty, then $\leg2p=1$ and every point satisfies
\begin{equation}\label{eq:common-ramified-classes}
 [\ell]_{\Q_p}\in\{1,[d]_{\Q_p}\}.
\end{equation}
If $\leg2p=1$, both classes in \eqref{eq:common-ramified-classes}
occur.
\end{lemma}
\begin{proof}
First suppose $e=v_p(s)>0$ is odd.  If $v_p(\ell)>0$, then
$P^2\equiv2\pmod p$.  Otherwise
$B^2=T(1+s/\ell)\equiv2\pmod p$.  Thus $2$ is a square modulo $p$.
A negative $n=v_p(\ell)$ makes $M^2=-s(1-T/\ell)$ have odd valuation $e$.
For $0<n<e$, the valuations $n$ of $C^2$ and $e-n$ of $M^2$
cannot both be even.  For $n>e$, $C^2$ has odd valuation $e$.
Only $n=0$ and $n=e$ remain.

If $n=0$, the third equation forces $v_p(T-\ell)>0$, so
$\ell\equiv T\equiv2\pmod p$ and $[\ell]_{\Q_p}=1$.
If $n=e$, then
\[
 M^2=(d/\ell)(1-\ell/T),
\]
and the second factor is a principal unit, hence a square.
Thus $[\ell]_{\Q_p}=[d]_{\Q_p}$.
The substitution
$(t,P,C,M,B)\mapsto(-t,C,P,B,M)$ leaves $\ell,d$ unchanged and proves
the assertion when $v_p(T)>0$.

For existence near $t=1$, choose $\ell=s+p^{2h}$ with $2h>e$.
Then $C=p^h$, and the other three radicands have square unit residues
$2,2,4$.  Hensel's lemma gives all three roots.  This realizes $[d]$;
the reflection \eqref{eq:common-reflection} realizes $1$.
The sign substitution treats the other case.
\end{proof}

\begin{lemma}[The dyadic exponent]\label{lem:common-dyadic}
If $\mathcal U_t(\Q_2)$ is nonempty, then either $v_2(t)\ge2$ or
$v_2(t)$ is negative and odd.  If additionally $d$ is a square in
$\Q_2$, then
\[
 v_2(\ell)=v_2(t)<0\quad\text{is odd}.
\]
\end{lemma}
\begin{proof}
Apply \eqref{eq:common-return}, without real sign restrictions, to obtain
coordinates $x,y$ and four distances $r_j$ in $\Q_2$.
Choose $N=2^h$ with $h\ge0$ minimal such that
$X=Nx,Y=Ny\in\Z_2$.  Then $A_j=Nr_j\in\Z_2$, since their squares
are integral.  If $h=0$, one of the four distances has square
$2\pmod4$, a contradiction.  Hence $h\ge1$, and $X,Y$ are not both
even.  They cannot both be odd, since $A_0^2=X^2+Y^2$.
After exchanging the coordinate axes, take $X$ odd and $Y$ even.
Every $A_j$ is odd.  The first and third squared distances modulo eight
now give $4\mid Y$ and $4\mid Y-N$, so $h\ge2$.

Put $U=N-2X$, $V=N-2Y$, and
$S_\pm=A_0A_3\pm A_1A_2$.  Direct expansion gives
\[
 S_+S_-=-N^2UV,\qquad t=S_+/N^2,\qquad v_2(U)=1.
\]
For $V\ne0$, we have $v_2(V)\ge2$.  Exactly one of the sum and
difference of two odd units has valuation one.  The product identity
therefore gives either $v_2(t)=1-2h$ or $v_2(t)=v_2(V)\ge2$.
If $V=0$, then $S_+\ne0$ because $t\ne0$, so $S_-=0$ and
$S_+=2A_0A_3$; again $v_2(t)=1-2h$.

If $d$ is a dyadic square, the alternative $v_2(t)\ge2$ is impossible:
it would give $d\equiv-1\pmod8$.  Write $m=v_2(t)<0$, which is odd,
and $n=v_2(\ell)$.  For $n>m$, $P^2=\ell+T$ has odd valuation $m$.
For $n<m$, $B^2=T(1+s/\ell)$ has odd valuation $m$.
Thus $n=m$.
\end{proof}

\subsection{The residue matrix}

Let $\Delta=p_1\cdots p_r>1$ be positive, odd and squarefree.  Define
$A_\Delta=(a_{ij})\in M_r(\F_2)$ by
\begin{equation}\label{eq:common-matrix}
 (-1)^{a_{ij}}=\leg{p_j}{p_i}\quad(i\ne j),\qquad
 a_{ii}=\sum_{j\ne i}a_{ij}.
\end{equation}
Thus every row sums to zero.  Quadratic reciprocity does not in general
make this matrix symmetric; the denominator $p_i$ in
\eqref{eq:common-matrix} specifies the place at which the row is imposed.

\begin{proposition}\label{prop:common-matrix}
Suppose
\begin{equation}\label{eq:common-delta-hyp}
 \Delta\equiv1\pmod8,\qquad p_i\equiv1\text{ or }7\pmod8,
 \qquad t>1,\qquad d\in\Delta\Q^{\times2}.
\end{equation}
Every rational point of $\mathcal U_t$ has
\begin{equation}\label{eq:common-kernel}
 [\ell]=\left[2\prod_{i=1}^r p_i^{e_i}\right],
 \qquad (e_1,\ldots,e_r)^\mathsf T\in\ker A_\Delta.
\end{equation}
If $\rank A_\Delta=r-1$, then
\begin{equation}\label{eq:common-two-classes}
 [\ell]\in\{[2],[2\Delta]\}.
\end{equation}
\end{proposition}
\begin{proof}
Write $d=\Delta q^2$.  An odd valuation of $d$ cannot be negative,
since $v_p(t)<0$ implies $v_p(t^2-1)=2v_p(t)$.
Thus $S(d)=\{p_1,\ldots,p_r\}$.
The support lemma leaves $[\ell]=[2^e\prod p_i^{e_i}]$.
As $\Delta\equiv1\pmod8$, $d$ is a dyadic square; the dyadic lemma
forces $e=1$.

At $p_i$, if $e_i=0$, the ramified-class lemma requires $\ell$ to be a
local square; if $e_i=1$, it requires $\ell/d$ to be a local square.
All valuation parities have already been accounted for.
Moreover $\leg2{p_i}=1$.  The remaining unit conditions are exactly
\[
 \sum_{j\ne i}a_{ij}(e_j+e_i)=0,
\]
which is row $i$ of \eqref{eq:common-kernel}.
The zero and all-one vectors belong to the kernel.  If its dimension
is one, they are its only vectors, proving \eqref{eq:common-two-classes}.
\end{proof}

\begin{lemma}[A local squareclass restriction at $37$]\label{lem:common-37}
Suppose $t\in\Z_{37}$ and
$\overline t\in\{14,-14,16,-16\}$.  Every point of
$\mathcal U_t(\Q_{37})$ has $[\ell]_{\Q_{37}}=1$.
\end{lemma}
\begin{proof}
The substitution $t\mapsto-t$ used above preserves $\ell$, so it is enough
to treat $\overline t=14,16$.  In either case $T$ is a nonsquare unit.
A positive valuation of $\ell$ would give $P^2\equiv T\pmod{37}$;
a negative valuation would give $B^2\equiv T\pmod{37}$.
Thus $\ell$ is a unit.  We prove that its reduction $\lambda$ is square.

Work in $\F_{37}$.  The elements $13,15,17$ are twice squares:
\[
 13=2\cdot5^2,\qquad15=2\cdot10^2,\qquad17=2\cdot8^2.
\]
The element $2$ is nonsquare, while $-1$ is square.
If $\lambda$ were nonsquare, the last two radicands, together with the
first two, would imply that all four elements
\[
 \lambda+\overline t+1,\quad\lambda+\overline t-1,
 \quad\lambda-\overline t+1,\quad\lambda-\overline t-1
\]
are squares, allowing zero.

Here is a polynomial test for two squares differing by $2$.
If $a^2=z+1$ and $b^2=z-1$, put $w=(a+b)^2/2$.
Since $(a+b)(a-b)=2$, we have $w\ne0$ and
\[
 z=(w+w^{-1})/2,\qquad w^{18}=-1.
\]
Let $T_j$ be the Chebyshev polynomials, characterized by
$T_j((w+w^{-1})/2)=(w^j+w^{-j})/2$.
Then $T_9(z)=0$.  Reduction of $-T_9/3$ modulo $37$ gives
\[
 F(Z)=Z^9+7Z^7+4Z^5+3Z^3-3Z.
\]
Apply this to $z=\lambda+\overline t$ and
$z=\lambda-\overline t$.  Both $F$-values vanish.  Expanding and
factoring their sum gives the identities in $\F_{37}[X]$
\begin{align*}
 F(X+14)+F(X-14)&=2X^3(X^2-14)(X^2+4)(X^2+6),\\
 F(X+16)+F(X-16)&=2X^3(X^2-14)(X^2+6)(X^2+18).
\end{align*}
The three elements $14,31=-6,19=-18$ are nonsquares: each has
$18$th power $-1$.  Since $\lambda\ne0$, the second identity is
impossible.  The first leaves only $\lambda^2=-4$, hence
$\lambda=12$ or $25$, both squares ($12=7^2$ and $25=5^2$).
This too contradicts the assumed nonsquareness.
Thus $\lambda$ is square; the square criterion for odd local fields
then gives $[\ell]_{\Q_{37}}=1$.
\end{proof}

\begin{theorem}[Residue-matrix exclusion]\label{thm:common-exclusion}\label{thm:common-squareclass}
Let $\Delta=p_1\cdots p_r>1$ be positive, odd and squarefree, and suppose
\[
 \Delta\equiv1\pmod8,\qquad
 p_i\equiv1\text{ or }7\pmod8,\qquad
 \rank_{\F_2}A_\Delta=r-1.
\]
For every rational $t>1$ satisfying
\begin{equation}\label{eq:common-excluded-t}
 t^2-1\in\Delta\Q^{\times2},\qquad t\in\Z_{37},\qquad
 \overline t\in\{14,-14,16,-16\},
\end{equation}
one has $D_t(\Q)=\varnothing$.
\end{theorem}
\begin{proof}
A rational point in the finite chart would satisfy
$[\ell]\in\{[2],[2\Delta]\}$ by Proposition~\ref{prop:common-matrix}.
The residues of $d$ at the two positive values of $\overline t$ are
$10=11^2$ and $33=12^2$.  The prime $37$ does not divide $\Delta$,
since its residue modulo eight is $5$.  From $d=\Delta q^2$ it follows
that $q$ is a $37$-adic unit and $\Delta$ is a square unit there.
Thus both $2$ and $2\Delta$ are nonsquares in $\Q_{37}$.
This contradicts Lemma~\ref{lem:common-37}.  Proposition~
\ref{prop:common-genus-boundary} excludes every rational point outside
the finite chart.
\end{proof}

The theorem constrains one global squareclass.  A different choice of
$\ell$ at each completion cannot replace the compatibility in
\eqref{eq:common-kernel}.  In particular, the theorem applies without a
bound on the height of $\ell$, on the height of $t$, or on the number of
prime factors of $\Delta$.

\subsection{An explicit infinite family}

Consider
\[
 \Delta_0=1926281441=7\cdot17\cdot23\cdot31\cdot73\cdot311.
\]
In the indicated prime order, the matrix \eqref{eq:common-matrix} is
\begin{equation}\label{eq:common-six-matrix}
 A_{\Delta_0}=\begin{pmatrix}
 0&1&0&1&1&1\\
 1&1&1&1&1&1\\
 1&1&0&0&0&0\\
 0&1&1&1&1&0\\
 1&1&0&1&1&0\\
 0&1&1&1&0&1
 \end{pmatrix}.
\end{equation}
Its entries follow from quadratic reciprocity and reduction of the six
primes modulo one another.  Its kernel has a short explicit
certificate.  Row three gives $e_1=e_2$; row five then gives
$e_4=e_5$; row four gives $e_2=e_3$; row six gives $e_4=e_6$;
and row one gives $e_2=e_4$.  Hence all six coordinates agree.
The rank is therefore five.

\begin{corollary}\label{cor:common-infinite}
Put $C_0=269297$ and $K=2^{20}\cdot7\cdot37$.  For every integer
$j\ge1$, define
\begin{equation}\label{eq:common-infinite}
 q_j=\frac1{C_0}+\frac K{1+K^2j},\qquad
 t_j=\frac{1+\Delta_0q_j^2}{1-\Delta_0q_j^2}.
\end{equation}
The parameters $t_j$ are distinct and satisfy
\[
 1<t_j<\frac{17}{16}<\frac{\sqrt5}{2},\qquad
 D_{t_j}(\Q)=\varnothing.
\]
Every $D_{t_j}$ has a real point in the strict return chamber
\eqref{eq:common-chamber}.  In addition, $t_j+1$ is neither a rational
square nor twice a rational square, and
$[t_j^2-1]=[\Delta_0]\notin\{1,[2]\}$.
\end{corollary}
\begin{proof}
For $q_0=1/C_0$, direct cancellation gives
\[
 t(q_0)=\frac{138225}{131072}=:t_0,\qquad
 1-\Delta_0q_0^2=\frac{2^{18}}{C_0},\qquad
 t_0\equiv14\pmod{37}.
\]
The perturbation $\delta_j=K/(1+K^2j)$ has $37$-adic valuation one.
Since the denominator of $t(q_0)$ is a $37$-adic unit,
$t_j\equiv14\pmod{37}$ for all $j$.  Also
\[
 t_j^2-1=\Delta_0
 \left(\frac{2q_j}{1-\Delta_0q_j^2}\right)^2.
\]
At the real place $0<\delta_j<1/K<10^{-7}$.  The integer inequality
\[
 33\Delta_0\left(\frac1{269297}+\frac1{10^7}\right)^2
 =\frac{24893364230896160241}{26929700000000000000}<1
\]
gives $0<\Delta_0q_j^2<1/33$, and hence $1<t_j<17/16$.
The sequence $q_j$ is strictly decreasing, and the function
$(1+\Delta_0q^2)/(1-\Delta_0q^2)$ is strictly increasing on
$0<q<1/\sqrt{\Delta_0}$.  Thus the $t_j$ are distinct.
Theorem~\ref{thm:common-exclusion} now applies to each one.

For the real assertion choose $\ell=\sqrt d$ and positive roots in
\eqref{eq:common-four-roots}.  Since $s<\sqrt{sT}<T$, all four
radicands are positive.  Furthermore
\[
 (PM)^2=(BC)^2=2\sqrt d.
\]
The inequality $t<\sqrt5/2$ gives $d<1/4$, so
$|PM|+|BC|=2\sqrt{2\sqrt d}<2$.
Finally $v_7(q_0)=-1$ and $v_7(\delta_j)=1$, whence
$v_7(q_j)=-1$ and $v_7(1-\Delta_0q_j^2)=-1$.
As $t_j+1=2/(1-\Delta_0q_j^2)$, we get $v_7(t_j+1)=1$,
which proves both square exclusions.
\end{proof}

The limit parameter itself also satisfies the exclusion theorem:
\begin{equation}\label{eq:common-t0}
 t_0=\frac{138225}{131072},\qquad
 t_0^2-1=\frac{\Delta_0}{2^{34}},\qquad D_{t_0}(\Q)=\varnothing.
\end{equation}
For a one-prime squareclass $\Delta=p\equiv1\pmod8$, the matrix is
$(0)$, so its rank hypothesis is automatic.  Thus
Theorem~\ref{thm:common-exclusion} also applies to every such $p$ and
every $t$ satisfying \eqref{eq:common-excluded-t}.

\subsection{The unramified genus-seventeen covers}

The support lemma also gives a uniform finite descent for each fixed
parameter.  For $t>1$ define
\[
 \mathcal K_t=\left\{2^e\prod_{p\in S(d)}p^{e_p}:
                         e,e_p\in\{0,1\}\right\}.
\]
For $\kappa\in\mathcal K_t$, let $\widetilde D_{t,\kappa}$ be the
smooth projective curve obtained by adjoining to
\eqref{eq:common-four-roots} the further equation
\begin{equation}\label{eq:common-genus17}
 Z^2=\ell/\kappa.
\end{equation}

\begin{proposition}\label{prop:common-genus17}
The morphism $\widetilde D_{t,\kappa}\to D_t$ is a geometrically
connected unramified double cover, and
$g(\widetilde D_{t,\kappa})=17$.  Every rational point of $D_t$ lifts
to one of these covers; on the affine chart the lift retains the same
coordinates $\ell,P,C,M,B$.
\end{proposition}
\begin{proof}
On $D_t$, the function $\ell$ has eight zeros and eight poles, all of order
two, since the map to the $\ell$-line has ramification index two over
$0$ and $\infty$.  Thus adjoining its square root is unramified.
To prove geometric connectedness, suppose that $\ell$ became a square in
the four-radical function field over $\overline\Q$.  In a
multiquadratic extension the base-field squareclasses that become
squares are exactly the span of its defining classes.
The valuations at $-T,s,T,-s$ force each coefficient of such an
expression for $[\ell]$ to vanish.  The remaining class $[\ell]$ is nonzero
on $\overline\Q(\ell)$, a contradiction.  Riemann--Hurwitz now gives
\[
 2g(\widetilde D_{t,\kappa})-2=2(2\cdot9-2)=32.
\]
All rational points of $D_t$ are in the affine chart by
Proposition~\ref{prop:common-genus-boundary}.  For such a point,
Lemma~\ref{lem:common-support} supplies $\kappa\in\mathcal K_t$ with
$[\ell]=[\kappa]$, giving a rational $Z$ in
\eqref{eq:common-genus17}.
\end{proof}

\begin{proposition}[Finite local compatibility]\label{prop:common-finite-places}
For fixed $t>1$, define
\[
 \mathfrak S_\ell(t)=\{\kappa\in\mathcal K_t:
    \widetilde D_{t,\kappa}(\Q_v)\ne\varnothing
    \text{ for every place }v\}.
\]
It suffices to impose these conditions at the real place, at the primes
$p\le1153$, and at the primes at which one of $t,t-1,t+1$ has nonzero
valuation.  If $\mathfrak S_\ell(t)$ is empty, then $D_t(\Q)$ is empty.
\end{proposition}
\begin{proof}
Let $p$ be an odd prime outside the stated set.  The parameters
$t,s,T,\kappa$ are units.  The six branch sections on the $\ell$-line
remain distinct modulo $p$: their differences are products of units
among $2,t,s,T$.  The normalization of the five-radical cover is
smooth over $\Z_p$.  More explicitly, away from the branch sections
it is finite \emph{\'etale}; near a branch section, after an
\'etale extension extracting square roots of unit factors, its local
models are $w^2=z$, with $z$ a parameter on the base.  These models
are smooth, and smoothness descends.  The private-valuation argument
in the preceding proposition persists on the geometric special fibre,
so the latter is connected of genus seventeen.

The Hasse--Weil bound gives
\[
 \#\widetilde D_{t,\kappa}(\F_p)\ge p+1-34\sqrt p>0
 \qquad(p\ge1163).
\]
The next prime after $1153$ is $1163$.  Smooth Hensel lifting therefore
supplies a $\Q_p$-point.  A smooth curve over a local field has
infinitely many points in a sufficiently small neighbourhood of any
rational point, so the point can also be chosen off the finite
boundary.  The final assertion follows by applying
Proposition~\ref{prop:common-genus17} to a rational point of $D_t$.
\end{proof}

Thus the common-coordinate descent produces finitely many
geometrically connected covers while retaining all four ordered roots.
The residue-matrix theorem supplies domains on which the associated
compatibility set is empty, and Corollary~\ref{cor:common-infinite}
gives an explicit infinite sequence of such domains inside the strict
real return range.
